\begingroup
\renewcommand{\thesection}{37C}
\section{Maxwell Evolution as Reciprocal Conservative Action}\label{sec:maxwelltranslation}
\status{conventional electromagnetic equations and explicit operator calculation; a native Thalean electromagnetic reduction remains open}
The useful translation is functional rather than letter-by-letter. Constraints restrict allowed field states; reciprocal derivatives specify their evolution; sources and boundaries exchange field energy; an operational information current additionally requires an encoding and a decoder. None of these roles is filled merely by naming six WXYZTI letters.

\subsection{Constraints, coupling, and continuity}
In SI vacuum units, with $c^2=(\epsilon_0\mu_0)^{-1}$, Maxwell's equations are \citep{feynmanmaxwell}
\begin{equation}
\begin{aligned}
\nabla\cdot\mathbf E&=\rho/\epsilon_0,&
\nabla\cdot\mathbf B&=0,\\
\partial_t\mathbf E&=c^2\nabla\times\mathbf B-\mathbf J/\epsilon_0,&
\partial_t\mathbf B&=-\nabla\times\mathbf E.
\end{aligned}
\tag{EM1}\label{eq:maxwelltranslation}
\end{equation}
The divergence equations serve as constraints on a time slice; describing them as registration constraints is a QR translation, not a claim that they are measurement operations. Taking divergences of the evolution equations gives
\begin{equation}
\partial_t(\nabla\cdot\mathbf B)=0,\qquad
\partial_t(\nabla\cdot\mathbf E-\rho/\epsilon_0)
=-\frac{\nabla\cdot\mathbf J+\partial_t\rho}{\epsilon_0}.
\tag{EM2}\label{eq:maxwellconstraints}
\end{equation}
Thus compatible initial constraints are preserved when charge continuity holds. The curl equations alone do not repair inconsistent initial constraints. Charge current $\mathbf J$, energy flux, and QR information current remain different quantities.

\subsection{A closed domain with an exact skew-adjoint generator}
Take constant vacuum coefficients and periodic fields on a flat three-torus $\mathbb T^3$. This is an explicit mathematical boundary choice, not a cosmological topology assertion. Put
\begin{equation}
\Psi=\begin{pmatrix}\mathbf E\\c\mathbf B\end{pmatrix},\qquad
\partial_t\Psi=\mathsf A_M\Psi,\qquad
\mathsf A_M=c\begin{pmatrix}0&\mathscr C\\-\mathscr C&0\end{pmatrix},
\quad \mathscr C=\nabla\times,
\quad \rho=0,\ \mathbf J=0.
\tag{EM3}\label{eq:maxwellgenerator}
\end{equation}
Use the $L^2$ inner product on the full six-component field. Periodic integration by parts makes curl symmetric; its Fourier realization on the domain of fields with square-integrable curl is self-adjoint. On that specified domain,
\[
\mathsf A_M^\dagger=-\mathsf A_M.
\]
Appendix~\ref{app:conservativecalculations} gives the domain and Fourier calculation, rather than inferring skew-adjointness from the appearance of a block matrix. On a bounded region, an adjoint calculation also involves its boundary domain; energy flux vanishing for one solution does not prove operator skew-adjointness \citep{ellerkarabash2021}.

The corresponding invariant is
\begin{equation}
\mathcal E_{\rm EM}
=\frac{\epsilon_0}{2}\int_{\mathbb T^3}
\bigl(|\mathbf E|^2+c^2|\mathbf B|^2\bigr)\,\mathrm d^3x
=\frac{\epsilon_0}{2}\|\Psi\|_{L^2}^2.
\tag{EM4}\label{eq:maxwellenergy}
\end{equation}
This is total electromagnetic energy in this conventional realization. It is not a pointwise invariant, a bit count, or a native QR complexity functional. Reversible source-free evolution of the represented field and energy invariance are both present here, but neither follows from the other in arbitrary systems.

\subsection{The local balance must remain visible}
With sources or a finite observation region $V$, the calculation instead gives Poynting's balance \citep{feynmanenergy}:
\begin{equation}
\begin{aligned}
u_{\rm EM}&=\tfrac12\bigl(\epsilon_0|\mathbf E|^2+|\mathbf B|^2/\mu_0\bigr),
&\mathbf S_P&=\mathbf E\times\mathbf B/\mu_0,\\
\partial_t u_{\rm EM}+\nabla\cdot\mathbf S_P&=-\mathbf J\cdot\mathbf E,\\
\frac{\mathrm d}{\mathrm dt}\mathcal E_V
&=-\int_{\partial V}\mathbf S_P\cdot\boldsymbol\nu\,\mathrm dA
-\int_V\mathbf J\cdot\mathbf E\,\mathrm d^3x.
\end{aligned}
\tag{EM5}\label{eq:poyntingbalance}
\end{equation}
Here $\boldsymbol\nu$ is the outward boundary normal. A field subsystem can gain or lose energy by boundary flux or work on charges. A closed field--matter model must include the matter, source apparatus, and any remaining boundary reservoir before asserting a conserved total energy. Prescribed sources alone do not supply that enlarged model or a proof of QR complete recoverability.

The analogy with propagated-plus-retained accounting concerns the need to state the complete boundary of the account. It does not identify matter energy with receipt entropy, or identify Poynting flux with operational information flux. A steady or standing field pattern need not encode a new distinguishable message merely because a local quadratic or cross-product expression is nonzero.

\subsection{Duality, signed reciprocity, and a second phase analogy}
On the source-free periodic domain, constant electric--magnetic duality rotations act by
\begin{equation}
\begin{aligned}
\Psi'&=e^{\vartheta\mathsf J}\Psi,
&\mathsf J&=\begin{pmatrix}0&I_3\\-I_3&0\end{pmatrix},
&\mathsf J^2&=-I_6,\\
\mathbf E'&=\mathbf E\cos\vartheta+c\mathbf B\sin\vartheta,
&c\mathbf B'&=-\mathbf E\sin\vartheta+c\mathbf B\cos\vartheta.
\end{aligned}
\tag{EM6}\label{eq:duality}
\end{equation}
The rotation preserves (EM4), the divergence-free subspace, and the evolution, since $[\mathsf J,\mathsf A_M]=0$. This is the conventional duality symmetry \citep{fernandezcorbaton2013}. A fixed electric source generally breaks it, and even a source-free boundary-value problem needs a duality-invariant boundary domain. Periodic boundaries satisfy that requirement; a perfectly conducting electric boundary generally does not.

The signed quarter-turn is not an unsigned reciprocal flip: $\mathsf J^2=-I_6$, whereas the swap $\mathsf P(E,cB)=(cB,E)$ has $\mathsf P^2=I_6$ and $\mathsf P\mathsf A_M\mathsf P=-\mathsf A_M$. The latter reverses the source-free evolution generator rather than commuting with it. Similarly, replacing the minus sign in (EM3) by a plus sign produces a generically nonconservative block operator. Reciprocal coupling therefore needs its signed law, not just two equal-sized sectors.

Duality rotates a field pair; the rotor fiber rotates a marked frame while fixing its axis. Equal $SO(2)$ group types do not identify their state spaces, readouts, or dynamics. In particular duality is not a gauge redundancy of the full electromagnetic field readout, and its angle is not the gyroscope's mechanical spin phase.

\subsection{Propagation adds more than skewness}
On the divergence-free subspace, squaring (EM3) yields the vacuum wave equations by $\nabla\times\nabla\times=\nabla\nabla\cdot-\nabla^2$. For a nonzero Fourier wavevector $k$, the symbol $\mathscr C(k)=i[k]_\times$ is Hermitian and the transverse branches have frequency $\omega=\pm c|k|$. These facts use the supplied Euclidean spatial derivative and Maxwell constitutive relation. Skewness alone allows many other dispersions and does not select Maxwell propagation, its polarizations, or its speed.

The physical $c$ in these equations is not a newly derived value of $c_{\rm Th}$. A QR reduction must still construct an admitted local field update and a calibration between its address, event, and physical scales. The common mathematical class established here is conservative skew evolution. A common native physical mechanism remains a separate theorem target.
\endgroup
\setcounter{section}{37}
